\documentclass[11pt,a4paper]{article}
\usepackage[utf8]{inputenc}
\usepackage{amsmath,amsthm,amssymb,mathtools}
\usepackage[margin=1in]{geometry}
\usepackage{hyperref}
\usepackage{cite}

\newtheorem{theorem}{Theorem}
\newtheorem{proposition}[theorem]{Proposition}
\newtheorem{lemma}[theorem]{Lemma}
\newtheorem{corollary}[theorem]{Corollary}
\newtheorem{conjecture}[theorem]{Conjecture}
\newtheorem{definition}[theorem]{Definition}

\title{Foundations of Event-Potential Theory:\\
Discrete general covariance and ballistic-time determination\\
of Hawkes kernels on causal graphs}

\author{[Author name]\\ [Affiliation]}
\date{\today}

\begin{document}
\maketitle

\begin{abstract}
We establish two foundational results for Event-Potential Theory (EPT), a framework that treats events as ontological primitives and the spacetime metric as emergent from event density on a causal graph. First, we prove that Hawkes processes on causal graphs satisfy discrete general covariance under three natural conditions: graph-distance kernel stationarity, causal locality, and vertex-homogeneous branching ratio. Second, we prove that the kernel itself is not a free function: under \emph{ballistic temporal concentration} of the Hawkes generation index (a structural theorem we establish via renewal theory), the kernel tail is determined by $(d-1)$-dimensional spatial path-counting, where $d$ is the spacetime dimension. The combined framework gives a one-parameter prediction for the Hurst exponent of macroscopic stochastic dynamics: $H_{\mathrm{EPT}}(d) = (d-4)/2$, with $d = 4$ (physical spacetime) sitting at the boundary of the rough Cox--Ingersoll--Ross class ($H = 0$). The boundary case is shown to correspond to a Gaussian multiplicative chaos (GMC) macroscopic limit with multifractal scaling. We verify the underlying tail prediction $\phi(\tau) \sim \tau^{-(d-1)/2}$ by simulation on $\mathbb{Z}^{1,d-1}$ for $d \in \{2,3,4\}$ with three-digit agreement to deterministic computation.
\end{abstract}

\section{Introduction}

The causal set programme \cite{BLMS1987} and related discrete approaches to quantum gravity confront a foundational challenge: identifying dynamics that respect background independence while producing macroscopic gravitational behavior. The Rideout--Sorkin classical sequential growth (CSG) family \cite{RS1999} provided an answer in the causal set setting, parametrizing covariant classical dynamics by a sequence of coupling constants.

This paper develops Event-Potential Theory (EPT) on parallel structural grounds. EPT takes events as primitive (rather than spacetime points), uses self-exciting point processes for dynamics, and treats the metric as emergent from event density. The natural dynamical engine for EPT is the Hawkes process \cite{Hawkes1971}, a continuous-time branching process whose self-excitation captures the idea that past events potentiate future ones.

We address two foundational questions:
\begin{enumerate}
\item \textbf{Is EPT structurally legitimate?} Does a Hawkes process on a causal graph satisfy discrete general covariance, the symmetry principle required of any background-independent discrete-spacetime framework?
\item \textbf{Is EPT predictively constrained?} Is the Hawkes kernel actually a free parameter, or is it determined by the combinatorial structure of the underlying graph plus structural symmetry requirements?
\end{enumerate}

We answer both affirmatively. The combined results establish EPT as a structurally legitimate (covariant) and predictively constrained (kernel-determined) framework with a single empirically-relevant free parameter $\eta$ (the branching ratio).

\section{Discrete general covariance for Hawkes-on-graph dynamics}

\subsection{Causal graphs and Hawkes processes}

A \emph{causal graph} is a triple $G = (V, E, \prec)$ where $V$ is a countable set of vertices, $E \subset V \times V$ is a set of directed edges, and $\prec$ is the transitive closure of $E$, satisfying antisymmetry. For $v \in V$, denote the past $J^-(v) = \{v' \prec v\}$ and future $J^+(v) = \{v \prec v'\}$.

The \emph{graph-causal distance} $d_G(v', v)$ for $v' \preceq v$ is the length of the longest directed path from $v'$ to $v$. This is a graph-isomorphism invariant.

A \emph{Hawkes process on $G$} is a family $\{N_v\}_{v \in V}$ of simple point processes with conditional intensities
\begin{equation}
\lambda_v(t) = \mu_v + \sum_{v' \preceq v} \int_{-\infty}^{t^-} \phi_{v',v}(t - s)\, dN_{v'}(s),
\end{equation}
where $\mu_v > 0$ is the baseline intensity at $v$ and $\phi_{v',v}: \mathbb{R}_+ \to \mathbb{R}_+$ is the kernel from $v'$ to $v$, vanishing unless $v' \preceq v$.

A \emph{labeled history} of duration $T$ is a sequence $H = ((v_1, t_1), \ldots, (v_N, t_N))$ with $t_1 \leq \cdots \leq t_N \leq T$. Two labeled histories are \emph{causally equivalent} if there exists a graph automorphism $\sigma$ of $G$ and a strictly increasing time reparametrization $\tau$ such that one history transforms into the other.

\subsection{The conditions}

\begin{description}
\item[(H1')] \textbf{Graph-distance stationarity.} $\phi_{v',v}(\tau) = \phi(\tau, d_G(v', v))$.
\item[(H2)] \textbf{Causal locality.} $\phi_{v',v}(\tau) = 0$ unless $v' \preceq v$.
\item[(H3)] \textbf{Vertex homogeneity.} Branching ratio $\eta_v = \eta$ and baseline $\mu_v = \mu$ independent of $v$.
\end{description}

\subsection{The covariance theorem}

\begin{theorem}[Discrete general covariance for Hawkes-on-graph]
\label{thm:dgc}
Let $G$ be a causal graph and let $\{N_v\}$ be a Hawkes process on $G$ satisfying (H1'), (H2), (H3). Let $P(H)$ denote the likelihood density of a labeled history $H$ of duration $T$. Then
\begin{equation}
P(H) = P(H')
\end{equation}
whenever $H$ and $H'$ are causally equivalent.
\end{theorem}

\begin{proof}
The Hawkes likelihood factors as
\begin{equation}
\mathcal{L}(H) = \prod_{i=1}^N \lambda_{v_i}(t_i^-) \cdot \exp\left(-\sum_v \int_0^T \lambda_v(t)\, dt\right).
\end{equation}
We show each factor depends only on the causal-equivalence class.

\textbf{Step 1: The exponential factor.} Under (H3), $\sum_v \int_0^T \lambda_v(t)\,dt = |V|\mu T + \sum_i \eta \int dt'$, depending only on $T$, $N$, and graph structure---all causet invariants.

\textbf{Step 2: The product factor.} For event $i$ at $(v_i, t_i)$,
\begin{equation}
\lambda_{v_i}(t_i^-) = \mu + \sum_{j: t_j < t_i, v_j \preceq v_i} \phi(t_i - t_j, d_G(v_j, v_i)).
\end{equation}
Under (H1'), each summand depends only on the time difference (invariant under simultaneous time-shift) and graph-causal distance (causet invariant). The summation set is the past light-cone of event $i$, itself a causet invariant.

Combining steps completes the proof.
\end{proof}

\subsection{CSG as projection of Hawkes-on-graph}

\begin{proposition}[Asymptotic CSG projection]
\label{prop:csg-projection}
Let $X$ be a Hawkes-on-graph process satisfying (H1'), (H2), (H3) with kernel $\phi(\tau, d)$. In the asymptotic regime where the kernel correlation time $\tau_{\mathrm{kernel}} \ll$ inter-event time $1/\Lambda$, the time-marginalized process is approximately a CSG dynamics with coupling constants
\begin{equation}
t_n^{\mathrm{eff}}(\phi) := \int_0^\infty \phi(\tau, n)\, d\tau.
\end{equation}
The approximation error is $O(\tau_{\mathrm{kernel}} \cdot \Lambda)$.
\end{proposition}

The proof and consequences are developed in Appendix~\ref{app:csg-proof}. The key point: CSG sits inside Hawkes-on-graph as the short-memory asymptotic skeleton. In the long-memory regime (which is precisely the near-critical regime relevant for emergent cosmology), Hawkes-on-graph extends rather than reduces to CSG.

\section{Ballistic time concentration in Hawkes branching}

\subsection{The generation index}

For a Hawkes process $N$ with kernel $\phi$ and branching ratio $\eta < 1$, the cluster representation \cite{HawkesOakes1974} organizes descendants into generations: an immigrant is generation 0, its direct offspring generation 1, and so on. For an event observed at time $\tau$ within a cluster rooted at time 0, let $G(\tau)$ denote its generation index.

The kernel has mean inter-generation interval $\delta_\phi = (1/\eta)\int_0^\infty \tau \phi(\tau)\,d\tau$.

\subsection{The ballistic concentration theorem}

\begin{theorem}[Ballistic concentration of generation index]
\label{thm:ballistic}
Let $N$ be a stationary Hawkes process with kernel $\phi$ satisfying:
\begin{enumerate}
\item[(K1)] $\eta < 1$ (subcriticality)
\item[(K2)] $\delta_\phi < \infty$ (finite first moment)
\item[(K3)] $\sigma_\phi^2 := (1/\eta)\int_0^\infty (s - \delta_\phi)^2 \phi(s)\,ds < \infty$ (finite second moment)
\end{enumerate}
For an event observed at time $\tau$ within its cluster (in the Palm sense), let $G(\tau)$ be its generation index. Then as $\tau \to \infty$:
\begin{enumerate}
\item[(a)] $G(\tau)/\tau \to 1/\delta_\phi$ in probability;
\item[(b)] $\displaystyle\frac{G(\tau) - \tau/\delta_\phi}{\sqrt{\tau \sigma_\phi^2/\delta_\phi^3}} \xrightarrow{d} N(0, 1)$;
\item[(c)] For any $\epsilon > 0$:
\begin{equation}
\mathbb{P}\left(|G(\tau) - \tau/\delta_\phi| > \epsilon\sqrt{\tau}\right) \leq \frac{1}{\epsilon^2} \cdot \frac{\sigma_\phi^2}{\delta_\phi^3} + o(1).
\end{equation}
\end{enumerate}
\end{theorem}

The proof combines: (i) the renewal CLT for iid sums \cite{Feller1971}, (ii) conditioned-walk Brownian-bridge convergence \cite{CsorgoHorvath1993}, (iii) Palm calculus for cluster processes \cite{DaleyVereJones2003}. Full proof in Appendix~\ref{app:ballistic-proof}.

\subsection{The structural consequence}

Theorem~\ref{thm:ballistic} states: the generation index at time $\tau$ is \emph{deterministically} $\tau/\delta_\phi + O(\sqrt{\tau})$, with relative fluctuation $O(1/\sqrt{\tau}) \to 0$. This is \emph{ballistic} behavior: deterministic linear growth, with diffusive fluctuations small relative to the mean.

The temporal direction in Hawkes branching is qualitatively different from spatial directions. \textbf{Time is ballistic; space is diffusive.}

\section{The Hawkes kernel from spatial path-counting}

\subsection{The reduction from $d$ to $d-1$ dimensions}

Consider a Hawkes process on the lattice $\mathbb{Z}^{1,d-1}$. The naive path-counting argument: paths of length $n$ have count $\sim (2d)^n/n^{d/2}$ by standard random-walk asymptotics, giving kernel tail $\phi(\tau) \sim \tau^{-d/2}$.

\emph{This argument fails.} It treats time as a diffusive lattice direction, but Theorem~\ref{thm:ballistic} establishes that time is ballistic: the generation index is sharply concentrated at $n \approx \tau/\delta_\phi$.

The correct argument: given Theorem~\ref{thm:ballistic} fixes the generation index deterministically, only the $(d-1)$ spatial directions contribute random-walk combinatorics:
\begin{equation}
P_n^{\mathrm{spatial}}(\vec{x}) \sim \frac{(2(d-1))^n}{n^{(d-1)/2}} \exp(-|\vec{x}|^2/(2n)).
\end{equation}
The $n^{(d-1)/2}$ denominator yields the corrected kernel tail.

\begin{theorem}[Corrected kernel tail]
\label{thm:tail}
Under (H1'), (H2), (H3), (K1)--(K3), the effective Hawkes propagator on $\mathbb{Z}^{1,d-1}$ has tail behavior
\begin{equation}
\phi_{\mathrm{eff}}(\tau) \sim \tau^{-(d-1)/2}
\end{equation}
as $\tau \to \infty$, up to logarithmic corrections at integer values of $(d-1)/2$.
\end{theorem}

The Jaisson--Rosenbaum kernel tail exponent (defined by $\phi(\tau) \sim \tau^{-(1+\alpha_{\mathrm{kernel}})}$) is therefore
\begin{equation}
\alpha_{\mathrm{kernel}}(d) = \frac{d-3}{2}.
\end{equation}
The corresponding Hurst exponent in the rough-CIR class \cite{JaissonRosenbaum2016} is
\begin{equation}
\boxed{H_{\mathrm{EPT}}(d) = \frac{d-4}{2}.}
\end{equation}

\subsection{Predictions across dimensions}

\begin{center}
\begin{tabular}{|c|c|c|l|}
\hline
$d$ & $\alpha_{\mathrm{kernel}}$ & $H$ & J-R class \\
\hline
2 & $-1/2$ & $-1$ & Below rough range \\
3 & $0$ & $-1/2$ & Below rough range \\
\textbf{4} & $\mathbf{1/2}$ & $\mathbf{0}$ & \textbf{Boundary of rough class} \\
5 & $1$ & $1/2$ & Boundary of standard class \\
6 & $3/2$ & $1$ & Light-tailed \\
\hline
\end{tabular}
\end{center}

Only $d=4$ (and marginally $d=5$) lie in the standard rough-CIR window $\alpha_{\mathrm{kernel}} \in (0,1)$; for $d \le 3$ the effective kernel is non-integrable and the Jaisson--Rosenbaum scaling limit does not apply in its standard form. This restriction is corroborated numerically in the companion simulation paper (Paper~3).

\section{Hawkes--wave correspondence}

\begin{proposition}[Hawkes--wave correspondence]
\label{prop:wave}
On the lattice $\mathbb{Z}^{1,d-1}$, the Hawkes propagator $\psi_\phi(\tau, v', v)$ and the discrete d'Alembertian Green's function $G_\Box^{\mathrm{ret}}(\tau, v', v)$ have the same asymptotic scaling form $\tau^{-(d-1)/2}$ times a Gaussian envelope in spatial separation. Specifically, for $\tau \to \infty$ with $|\vec{x}| = O(\sqrt{\tau})$:
\begin{equation}
\frac{\psi_\phi(\tau, \vec{x})}{G_\Box^{\mathrm{ret}}(\tau, \vec{x})} \to C_\phi
\end{equation}
where $C_\phi$ is a kernel-dependent constant.
\end{proposition}

The proof combines ballistic time concentration (Theorem~\ref{thm:ballistic}) with the local CLT for $(d-1)$-dimensional spatial random walks. Exact matching including all constants and subleading corrections requires a more refined combinatorial analysis (e.g., via Lindström--Gessel--Viennot path-sum determinants \cite{Lindstrom1973,GesselViennot1985}) and is left to follow-up work.

The structural message: Hawkes-on-graph dynamics with our conditions are not arbitrary self-excitation processes but natural realizations of discrete wave propagation: ballistic in time, diffusive in space.

\section{Numerical confirmation}

We tested the prediction $\phi_{\mathrm{eff}}(\tau) \sim \tau^{-(d-1)/2}$ via simulation of Hawkes processes on $\mathbb{Z}^{1,d-1}$ for $d \in \{2, 3, 4\}$. Methodology: Hawkes-on-lattice simulation via Ogata thinning with exponential elementary kernels; measurement of expected intensity propagator $\psi(\tau, d)$ from seed events; tail exponent extraction via Hill estimator, maximum-likelihood power-law fit, and an exact deterministic lattice computation (random-walk return probabilities convolved with Erlang generation weights). The full methodology, code, and figures are reported in the companion simulation paper (Paper~3).

\begin{center}
\begin{tabular}{|c|c|c|c|}
\hline
$d$ & Predicted $(d-1)/2$ & Exact computation & MC fit ($\eta=0.99$, 95\% CI) \\
\hline
2 & 0.500 & 0.4975 & $0.587\ [0.555, 0.622]$ \\
3 & 1.000 & 0.9970 & $1.051\ [1.014, 1.094]$ \\
4 & 1.500 & 1.4985 & $1.533\ [1.484, 1.583]$ \\
\hline
\end{tabular}
\end{center}

All three methods (exact deterministic computation, Monte-Carlo with block-bootstrap CIs, truncated power-law MLE) land on $(d-1)/2$ and exclude the naive $d/2$ prediction. The exact computation pins each value to three significant figures independent of statistical noise. The Monte-Carlo fits sit modestly above $(d-1)/2$ owing to finite-$\eta$ cutoff and finite fit-window bias, but every confidence interval excludes $d/2$. The structural prediction is confirmed.

\section{The $d = 4$ boundary case: GMC macroscopic limit}

\subsection{Spectral mechanism for log-correlation}

The $d = 4$ prediction $\alpha_{\mathrm{kernel}} = 1/2$ corresponds to kernel tail $\phi(\tau) \sim C\tau^{-3/2}$ (with appropriate UV regularization). The Fourier transform satisfies at small $\omega$:
\begin{equation}
1 - \hat\phi(\omega) \approx C_1\sqrt{i\omega \tau_0}
\end{equation}
where $C_1$ depends on UV regularization and $\tau_0$ is the UV cutoff.

\begin{lemma}[Log-correlated autocovariance at $d=4$]
\label{lem:log-cor}
For nearly-unstable Hawkes with kernel tail $\phi(\tau) \sim C\tau^{-3/2}$:
\begin{enumerate}
\item[(i)] $\hat\phi(\omega)$ has square-root branch structure at $\omega = 0$.
\item[(ii)] At criticality ($\eta \to 1$), the Bartlett spectral density satisfies $\hat\nu(\omega) \sim \Lambda/(C_1^2|\omega|\tau_0)$ at intermediate $\omega$.
\item[(iii)] The autocovariance has form
\begin{equation}
\nu(\tau) \sim A - K\log|\tau| + \text{finite}
\end{equation}
in the intermediate range $\tau_0 \ll \tau \ll T_c/(1-\eta)$, with explicit coefficient $K = \Lambda/(C_1^2\tau_0)$.
\item[(iv)] The variance of the integrated intensity scales as $O(T \log T)$ in the same intermediate range.
\end{enumerate}
\end{lemma}

\begin{proof}
(i) follows from $\int_0^\infty \tau^{-3/2} e^{-(1/T_c + i\omega)\tau}d\tau = \Gamma(-1/2)(1/T_c + i\omega)^{1/2}$. (ii) follows from $|1-\hat\phi(\omega)|^2 \approx C_1^2|\omega|\tau_0$ near criticality. (iii) is the inverse Fourier transform of $1/|\omega|$, which has logarithmic structure. (iv) follows from integration of $\nu(\tau)$ over $[0, T]$ with appropriate cutoffs.
\end{proof}

\subsection{The GMC macroscopic limit}

The log-correlated covariance structure is the foundational input to Gaussian multiplicative chaos (GMC) theory \cite{RhodesVargas2014,Berestycki2017}. Define the renormalized exponential measure
\begin{equation}
M_\epsilon(dt) := \exp\left(\gamma Y_\epsilon(t) - \frac{\gamma^2}{2}\mathbb{E}[Y_\epsilon(t)^2]\right) dt
\end{equation}
where $Y_\epsilon$ is the smoothed centered intensity at scale $\epsilon$.

\begin{theorem}[GMC macroscopic limit at $d=4$]
\label{thm:gmc}
For nearly-unstable Hawkes satisfying the $d = 4$ EPT structural prediction (kernel tail $\tau^{-3/2}$), and for $\gamma \in (0, \sqrt{2})$, the renormalized measure $M_\epsilon$ converges weakly as $\epsilon \to 0$ to a Gaussian multiplicative chaos measure $M_\gamma$ with logarithmic covariance kernel determined by the EPT parameters. The limit $M_\gamma$ has:
\begin{enumerate}
\item[(i)] Multifractal scaling: $\mathbb{E}[M_\gamma([0,T])^q] \sim T^{q + q(1-q)\gamma^2/(2K)}$;
\item[(ii)] Random measure structure (non-atomic with full support);
\item[(iii)] Non-trivial total mass for $\gamma < \sqrt{2}$.
\end{enumerate}
\end{theorem}

The proof applies Berestycki's GMC convergence criteria \cite{Berestycki2017} to the log-correlated structure established in Lemma~\ref{lem:log-cor}. Verification of the conditions is given in Appendix~\ref{app:gmc-conditions}.

\subsection{Empirical signature}

The macroscopic stochastic intensity (or $\Lambda(t)$ in cosmological applications) has multifractal scaling distinguishing this case from:
\begin{itemize}
\item Standard CIR / Sorkin Poissonian $\Lambda$ ($H = 1/2$, monofractal)
\item Rough CIR with $H > 0$ (monofractal at $H$)
\item $\Lambda$CDM (deterministic constant)
\end{itemize}
The multifractal spectrum $\xi(q) = q + q(1-q)\gamma^2/(2K)$ is the cleanest empirical distinction. Its cosmological consequences are developed in Paper~2.

\section{The critical role of $d = 4$}

That physical spacetime $d = 4$ sits exactly at the boundary of the rough CIR class ($H = 0$, log-correlated limit) is structurally striking. Three possible interpretations:
\begin{itemize}
\item \textbf{Coincidence.} The values are combinatorially forced; their match to physical $d = 4$ has no deeper meaning.
\item \textbf{Dynamical selection.} Causal graphs supporting stable spacetime structure are constrained to universality boundaries.
\item \textbf{Anthropic.} Observers can exist only in spacetimes near criticality.
\end{itemize}
We refrain from committing to an interpretation but flag the structural fact: the dimension producing critical-roughness behavior in the Hawkes macroscopic limit coincides with the observed spacetime dimension.

\section{Conclusions and open questions}

We have established two foundational results: (1) Hawkes-on-graph dynamics satisfy discrete general covariance (Theorem~\ref{thm:dgc}); (2) the kernel is determined by spatial path-counting via ballistic time concentration (Theorems~\ref{thm:ballistic}, \ref{thm:tail}, Proposition~\ref{prop:wave}). The combined framework gives a one-parameter prediction $H_{\mathrm{EPT}}(d) = (d-4)/2$ with $d = 4$ producing a GMC macroscopic limit (Theorem~\ref{thm:gmc}).

Open questions:
\begin{enumerate}
\item[(A)] Full proof of exact $d$'Alembertian matching via LGV path-sum determinants.
\item[(B)] Extension to non-lattice (Poisson-sprinkled) causal graphs.
\item[(C)] Computation of multifractal exponents in terms of EPT parameters $(\eta, T_c, \mu)$.
\item[(D)] Quantum extension via Hudson--Parthasarathy stochastic calculus.
\end{enumerate}

\appendix

\section{Proof of Proposition~\ref{prop:csg-projection} (sketch)}
\label{app:csg-proof}

Time-marginalization analysis: show the probability of the next vertex factors through integrated kernel norms in the short-memory regime $\tau_{\mathrm{kernel}} \cdot \Lambda \ll 1$. Standard point-process integration techniques apply. The full calculation is approximately two pages of explicit work and is available on request.

\section{Proof of Theorem~\ref{thm:ballistic} (full)}
\label{app:ballistic-proof}

\subsection{Renewal CLT for single ancestral chain}

For iid waiting times $W_1, W_2, \ldots$ with common density $h = \phi/\eta$, mean $\delta_\phi$, variance $\sigma_\phi^2$, the renewal counting function $N(\tau) = \max\{n: T_n \leq \tau\}$ satisfies (Feller \cite{Feller1971} §XI.5):
\begin{equation}
\frac{N(\tau) - \tau/\delta_\phi}{\sqrt{\tau\sigma_\phi^2/\delta_\phi^3}} \xrightarrow{d} N(0,1).
\end{equation}

\subsection{Lemma B.1: Size-biased waiting times}

\begin{lemma}
For $(W_1, \ldots, W_n)$ iid with density $h$, conditional on $T_n \in [\tau, \tau+\varepsilon]$ with $n = \lfloor\tau/\delta_\phi\rfloor$, as $\tau \to \infty$:
\begin{enumerate}
\item[(i)] The marginal density of any $W_i$ converges to $h$ pointwise with rate $O(1/n)$.
\item[(ii)] The empirical process $(T_{\lfloor nt\rfloor} - \lfloor nt\rfloor\delta_\phi)/\sqrt{n}$ converges to a Brownian bridge on $[0,1]$ \cite{CsorgoHorvath1993}.
\item[(iii)] The renewal CLT holds uniformly: $(N(\tau') - \tau'/\delta_\phi)/\sqrt{\tau'\sigma_\phi^2/\delta_\phi^3} \to N(0,1)$ for $\tau' = c\tau$, $c \in (0,1)$.
\end{enumerate}
\end{lemma}

\begin{proof}
By exchangeability and Bayes:
\begin{equation}
q_n(w; \tau) = h(w) \cdot \frac{p_{n-1}(\tau - w)}{p_n(\tau)}
\end{equation}
where $p_k$ is the $k$-fold convolution of $h$. By the local CLT \cite{Petrov1975} for $p_n$ at $n = \lfloor\tau/\delta_\phi\rfloor$:
\begin{equation}
p_n(\tau) \approx \frac{1}{\sqrt{2\pi n\sigma_\phi^2}}.
\end{equation}
The ratio $p_{n-1}(\tau-w)/p_n(\tau) = \sqrt{n/(n-1)}(1 + O(1/n)) \to 1$, establishing (i) with explicit rate. Statements (ii) and (iii) follow from the Donsker-type theorem for conditioned random walks (Csörgő-Horváth Theorem 3.2.1) and continuous mapping.
\end{proof}

\subsection{Lemma B.2: Palm distribution analysis}

\begin{lemma}
For the stationary Hawkes cluster process $X$, the Palm-Khinchin theorem \cite{DaleyVereJones2003} together with the Slivnyak-Mecke theorem for Poisson cluster processes gives, for an event observed at the origin under the Palm distribution, the generation $G(\tau)$ at separation $\tau$ from the cluster root satisfies the CLT scaling of Lemma B.1.
\end{lemma}

\begin{proof}
The Palm-Khinchin theorem decomposes the Palm distribution at the origin into immigrant plus offspring contributions. The Slivnyak-Mecke property of Poisson cluster processes gives: conditional on an event at origin, an independent additional cluster rooted at origin plus an event of an existing cluster with seed time $T_{\mathrm{seed}}$. The latter has distribution concentrated at $T_{\mathrm{seed}} \approx -\tau$ with width $O(\sqrt{\tau})$, and Lemma B.1(iii) gives uniform CLT scaling over this range.
\end{proof}

\subsection{Combining}

Theorem~\ref{thm:ballistic} follows from Lemmas B.1 and B.2 by combining the renewal CLT for the ancestor chain with the Palm-distribution decomposition. Statement (b) is the immediate combination; (a) and (c) follow by Slutsky's theorem and Chebyshev's inequality.

\section{GMC convergence conditions for Theorem~\ref{thm:gmc}}
\label{app:gmc-conditions}

Berestycki's elementary approach to GMC \cite{Berestycki2017} requires:
\begin{itemize}
\item[(B1)] Logarithmic covariance: $K_\epsilon(s,t) \to -\log|s-t| + g(s,t)$ with $g$ bounded continuous.
\item[(B2)] Uniform bounds: $K_\epsilon(s,s) \leq C - \log\epsilon$ and $K_\epsilon(s,t) \leq C - \log(|s-t| \vee \epsilon)$.
\item[(B3)] Convergence rate: uniform on compact sets away from diagonal.
\end{itemize}

For our smoothed Hawkes intensity at $d = 4$, Lemma~\ref{lem:log-cor} establishes (B1) with $g = A/K$ constant. The smoothing prescription $Y_\epsilon(t) = (\sqrt{K})^{-1}\epsilon^{-1}\int_{t-\epsilon/2}^{t+\epsilon/2}\delta\lambda(s)\,ds$ gives bounded $K_\epsilon(t,t) \sim -\log\epsilon$ at $\epsilon$ small, satisfying (B2). Uniform convergence (B3) follows from spectral continuity of the Hawkes process. Conditions verified; Berestycki's theorem applies.

\begin{thebibliography}{99}

\bibitem{BLMS1987} Bombelli, L., Lee, J., Meyer, D., Sorkin, R.D. (1987). Spacetime as a causal set. \emph{Phys. Rev. Lett.} \textbf{59}, 521.

\bibitem{RS1999} Rideout, D.P., Sorkin, R.D. (1999). A classical sequential growth dynamics for causal sets. \emph{Phys. Rev. D} \textbf{61}, 024002. arXiv:gr-qc/9904062.

\bibitem{Hawkes1971} Hawkes, A.G. (1971). Spectra of some self-exciting and mutually exciting point processes. \emph{Biometrika} \textbf{58}, 83.

\bibitem{HawkesOakes1974} Hawkes, A.G., Oakes, D. (1974). A cluster process representation of a self-exciting process. \emph{J. Appl. Probab.} \textbf{11}, 493.

\bibitem{Feller1971} Feller, W. (1971). \emph{An Introduction to Probability Theory and Its Applications, Vol. II}, 2nd ed., Wiley.

\bibitem{Petrov1975} Petrov, V.V. (1975). \emph{Sums of Independent Random Variables}, Springer.

\bibitem{CsorgoHorvath1993} Csörgő, M., Horváth, L. (1993). \emph{Weighted Approximations in Probability and Statistics}, Wiley.

\bibitem{DaleyVereJones2003} Daley, D.J., Vere-Jones, D. (2003). \emph{An Introduction to the Theory of Point Processes, Vol. I and II}, Springer.

\bibitem{JaissonRosenbaum2015} Jaisson, T., Rosenbaum, M. (2015). Limit theorems for nearly unstable Hawkes processes. \emph{Ann. Appl. Probab.} \textbf{25}, 600.

\bibitem{JaissonRosenbaum2016} Jaisson, T., Rosenbaum, M. (2016). Rough fractional diffusions as scaling limits of nearly unstable heavy tailed Hawkes processes. \emph{Ann. Appl. Probab.} \textbf{26}, 2860.

\bibitem{BenincasaDowker2010} Benincasa, D.M.T., Dowker, F. (2010). The scalar curvature of a causal set. \emph{Phys. Rev. Lett.} \textbf{104}, 181301.

\bibitem{Lindstrom1973} Lindström, B. (1973). On the vector representations of induced matroids. \emph{Bull. London Math. Soc.} \textbf{5}, 85.

\bibitem{GesselViennot1985} Gessel, I., Viennot, G. (1985). Binomial determinants, paths, and hook length formulae. \emph{Adv. Math.} \textbf{58}, 300.

\bibitem{RhodesVargas2014} Rhodes, R., Vargas, V. (2014). Gaussian multiplicative chaos and applications: A review. \emph{Probab. Surveys} \textbf{11}, 315.

\bibitem{Berestycki2017} Berestycki, N. (2017). An elementary approach to Gaussian multiplicative chaos. \emph{Electron. Commun. Probab.} \textbf{22}, 27.

\end{thebibliography}

\end{document}
